% Generated draft: requires mathematical and source review.
\documentclass[11pt]{article}
\usepackage[margin=25mm]{geometry}
\usepackage{fontspec}
\setmainfont{Latin Modern Roman}
\usepackage{amsmath,amssymb,mathtools,unicode-math}
\setmathfont{Latin Modern Math}
\usepackage{enumitem,needspace,xurl,hyperref}
\setlength{\parindent}{0pt}
\setlength{\parskip}{5pt}
\setlength{\emergencystretch}{3em}
\newcommand{\N}{\mathbb{N}}
\newcommand{\Z}{\mathbb{Z}}
\newcommand{\Q}{\mathbb{Q}}
\newcommand{\R}{\mathbb{R}}
\newcommand{\C}{\mathbb{C}}
\newcommand{\F}{\mathbb{F}}
\newcommand{\A}{\mathbb{A}}
\newcommand{\PP}{\mathbb{P}}
\newcommand{\E}{\mathbb{E}}
\newcommand{\eps}{\varepsilon}
\newcommand{\dd}{\,\mathrm{d}}
\begin{document}
\section*{UnsolvedMath ID: 20003337}
\textbf{Source problem code:} catalogue entry
\subsection{Definitions and mathematical statement}
J. Marklof's Problem 44 asks to classify quantum limits of eigenfunctions of a polygonal billiard. The supplied verification fixes Dirichlet \emph{configuration-space} limits: for $L^2$-normalized Dirichlet eigenfunctions $\psi_j$ on a bounded polygon $P$ with eigenvalues tending to infinity, these are weak limits of $|\psi_j|^2\,dx$. A phase-space semiclassical measure is a different object. The general polygon class, degeneracy choices and the full classification are not specified by the one-sentence source.

The aid treats rational polygons with finite reflection group $G$ and compact translation-surface unfolding $X(P)$, formed from $|G|$ copies of $P$. Let $q:X(P)\to P$ be folding and let the Laplacian on the unfolding be the Friedrichs extension at conical points. With $\chi(g)=\det g$, the reported unitary correspondence is
\[
 (U_\chi\psi)(g,p)=|G|^{-1/2}\chi(g)\psi(p).
\]
It identifies Dirichlet eigenfunctions on $P$ with the determinant-character sector of the unfolding. The analogous formula with $\chi=1$ identifies Neumann eigenfunctions with the trivial sector. The normalization gives equal $L^2$ norms because there are $|G|$ isometric copies. Boundary traces and conical extension conventions must be matched; no arbitrary self-adjoint extension is substituted.

Absolute squares in either sector are $G$-invariant. The aid reports a bijection between billiard configuration quantum limits and \emph{sector} quantum limits on $X(P)$, by folding and the unique $G$-invariant probability lift. The lift can be described by averaging a test function over the finite $G$-orbit of any lift of a point of $P$, with stabilizers counted through that average. It is not a claim that every $G$-invariant quantum limit produced by arbitrary unfolding eigenfunctions is necessarily realized by the chosen sector.

For regular phase-space limits the summary reports a corresponding invariant quotient dictionary away from cone/corner and glancing sets. Control of possible measure on those exceptional sets is not supplied, so no unrestricted global phase-space correspondence is asserted. The configuration result does not automatically resolve those additional issues.

If the unfolding is a flat two-torus, the aid invokes Jakobson's theorem that every toral eigenfunction configuration limit is absolutely continuous. Finite-piece folding then gives the same conclusion for this billiard symmetry sector, excluding support on a single periodic billiard trajectory of zero planar area. It does not exclude atomic momentum directions or transfer the conclusion to every higher-genus translation unfolding or irrational polygon. This is a restricted obstruction, not a full classification of polygonal quantum limits.

The appended Section 6.4 is separate Berry--Tabor context. It says generic integrable Laplace spectra should have Poisson statistical behavior [133] and describes two-torus eigenvalues as $\{Q(x):x\in\Z^2\}$ for positive-definite $Q$. Statistical scaling, multiplicity/symmetry conventions and the meaning of generic are not specified there. Those conjectural spectral-statistics assertions are not clauses or proven consequences of Problem 44.
\subsection{Short English statement}
The aid gives a boundary-character-sensitive unfolding reduction for rational polygons, preserving the exact normalization. Torus unfoldings then inherit a reported absolute-continuity obstruction for configuration limits. This does not classify all polygonal or phase-space limits, and the following Berry–Tabor paragraph is separate context.
\subsection{Sources}
\begin{itemize}
\item Local UnsolvedMath registry, permanent ID 20003337, source code catalogue entry; \texttt{aim-other-notes.json}, notes[143].
\item AIM, \emph{Emerging applications of measure rigidity}, Problem 44, J. Marklof, followed by Section 6.4: \url{https://aimath.org/WWN/measrigid/measrigid.pdf}. The supplied verification separates the one-sentence problem from that subsection; this draft does not independently check it.
\item Jakobson's toral absolute-continuity theorem and the unfolding/Friedrichs-sector dictionary are invoked by the aid without complete bibliographic entries or theorem statements.
\item Neighboring Berry--Tabor conjecture reference [133] is supplied without a full entry.
\item Machine-generated research attempt 1: partial result, open at the research run, rigor labelled ``proved,'' low novelty confidence; registry partially solved. Not peer reviewed; independently verify before citation. Registry updated 14 August 2026.
\end{itemize}
\end{document}
